

\medskip
\noindent\SGAInlineRefTarget{sga3:iii:corollary:4-31}\textbf{Corollary 4.31.}
\textit{Let \(Y\) be as in \SGARefLink{sga3:iii:lemma:4-27}{4.27}. The set of conjugates of \(Y\) by
sections of \(X\) over \(S\) inducing the unit section of
\(X_{\mathscr J}\) is isomorphic to}
\[
 E=\partial C^0(Y_0,L_0^X/L_0^Y)
 =C^0(Y_0,L_0^X/L_0^Y)\big/H^0(Y_0,L_0^X/L_0^Y).
\]

\medskip
\noindent\SGAInlineRefTarget{sga3:iii:corollary:4-32}\textbf{Corollary 4.32.}
\textit{Suppose in addition that \(S_0\) is affine and that
\(\omega^1_{Y_0/S_0}\) is finite locally free. If}
\[
 \mathscr F_0=
 \bigl(\operatorname{Lie}(X_0/S_0)/
       \operatorname{Lie}(Y_0/S_0)\bigr)
 \otimes_{\mathscr O_{S_0}}\mathscr J,
\]
\textit{then
\(E=\Gamma(S_0,\mathscr F_0)/H^0(Y_0,\mathscr F_0)\).}

Indeed,\footnote{Editor's note: The original has been expanded in what
follows.} since \(\omega^1_{Y_0/S_0}\) is finite locally free, as is
\(\omega^1_{X_0/S_0}\) (because \(X\) is assumed smooth over \(S\)),
\SGARefLink{sga3:iii:remark:0-6}{0.6} gives
\[
 L_0^Y=\mathbf W\bigl(
   \operatorname{Lie}(Y_0/S_0)\otimes_{\mathscr O_{S_0}}\mathscr J
 \bigr),
 \qquad
 L_0^X=\mathbf W\bigl(
   \operatorname{Lie}(X_0/S_0)\otimes_{\mathscr O_{S_0}}\mathscr J
 \bigr).
\]
On the other hand, by \SGARefLink{sga3:iii:lemma:4-25}{4.25} there is an exact sequence of
\(Y_0\)-\(\mathscr O_{S_0}\)-modules
\[
 0\longrightarrow \mathscr K\longrightarrow \omega^1_{X_0/S_0}
 \xrightarrow{\varphi}\omega^1_{Y_0/S_0}\longrightarrow0,
 \qquad \mathscr K=\ker(\varphi).
\]
Since \(\omega^1_{Y_0/S_0}\) and \(\omega^1_{X_0/S_0}\) are finite
locally free, there is a locally split exact sequence
\[
 0\longrightarrow
 \operatorname{Lie}(Y_0/S_0)\otimes_{\mathscr O_{S_0}}\mathscr J
 \longrightarrow
 \operatorname{Lie}(X_0/S_0)\otimes_{\mathscr O_{S_0}}\mathscr J
 \longrightarrow \mathscr F_0\longrightarrow0.
\]
It follows that there is an exact sequence of
\(Y_0\)-\(\mathscr O_{S_0}\)-modules
\[
 0\longrightarrow L_0^Y\longrightarrow L_0^X
 \longrightarrow\mathbf W(\mathscr F_0).
\]
By the argument used to prove \SGARefLink{sga3:iii:corollary:4-31}{4.31}, we may compute \(E\) as the image
of the composite map
\[
 C^0(Y_0,L_0^X)\xrightarrow{\pi}
 C^0(Y_0,\mathbf W(\mathscr F_0))\xrightarrow{\partial}
 C^1(Y_0,\mathbf W(\mathscr F_0)).
\]
Now \(\pi\) is the map
\[
 \Gamma\bigl(S_0,
   \operatorname{Lie}(X_0/S_0)\otimes_{\mathscr O_{S_0}}\mathscr J
 \bigr)
 \longrightarrow\Gamma(S_0,\mathscr F_0).
\]
Since \(S_0\) is affine, \(\pi\) is surjective, and this proves the
assertion.

\medskip
\noindent\SGAInlineRefTarget{sga3:iii:corollary:4-33}\textbf{Corollary 4.33.}
\textit{Let \(S,\mathscr I,\mathscr J\), and \(X\) be as in \SGARefLink{sga3:iii:theorem:4-21}{4.21},
and let \(Y\) be a diagonalizable subgroup of \(X\). Suppose that
\(\omega^1_{Y_0/S_0}\) is finite locally free and that \(S_0\) is
affine. The set of
subgroups of \(X\) conjugate to \(Y\) by a section of \(X\) over
\(S\) inducing the unit section of \(X_{\mathscr J}\) is isomorphic
to}
\[
 E=\Gamma\!\left(
 S_0,
 \frac{\operatorname{Lie}(X_0/S_0)}
      {\operatorname{Lie}(X_0/S_0)^{\operatorname{ad}(Y_0)}}
 \otimes_{\Gamma(S_0,\mathscr O_{S_0})}\Gamma(S_0,\mathscr J)
 \right),
\]
\textit{that is, to
\(L_0^X(Y_0)/H^0(Y_0,L_0^X)\).}\footnote{Editor's note: What follows
has been added; cf. \SGARefLink{sga3:iii:corollary:4-34}{4.34}.}

Indeed, by I, \SGARefLink{sga3:i:proposition:4-7-3}{4.7.3} (cf. \SGARefLink{sga3:iii:corollary:2-13}{2.13}), write
\[
 \operatorname{Lie}(X_0/S_0)
 =\operatorname{Lie}(X_0/S_0)^{\operatorname{ad}(Y_0)}
  \oplus\mathscr R.
\]
Since \(Y_0\) is commutative,
\[
 \operatorname{Lie}(Y_0/S_0)
 \subset\operatorname{Lie}(X_0/S_0)^{\operatorname{ad}(Y_0)},
\]
and hence
\[
 \begin{aligned}
 \mathscr F_0={}&
 \left(
  \frac{\operatorname{Lie}(X_0/S_0)^{\operatorname{ad}(Y_0)}}
       {\operatorname{Lie}(Y_0/S_0)}
 \right)\otimes_{\mathscr O_{S_0}}\mathscr J
 \ \oplus\ \mathscr R\otimes_{\mathscr O_{S_0}}\mathscr J,\\
 \mathscr F_0^{\operatorname{ad}(Y_0)}={}&
 \left(
  \frac{\operatorname{Lie}(X_0/S_0)^{\operatorname{ad}(Y_0)}}
       {\operatorname{Lie}(Y_0/S_0)}
 \right)\otimes_{\mathscr O_{S_0}}\mathscr J.
 \end{aligned}
\]
By \SGARefLink{sga3:iii:corollary:4-32}{4.32}, therefore,
\(E\simeq
\Gamma(S_0,\mathscr R\otimes_{\mathscr O_{S_0}}\mathscr J)\).
Returning to the definition of \(\mathscr R\) completes the proof.

\medskip
\noindent\SGAInlineRefTarget{sga3:iii:corollary:4-34}\textbf{Corollary 4.34.}
\textit{Let \(S,\mathscr I,\mathscr J\), and \(X\) be as in \SGARefLink{sga3:iii:theorem:4-21}{4.21},
and let \(Y\) be a diagonalizable subgroup of \(X\). Suppose that
\(\omega^1_{Y_0/S_0}\) is finite locally free and that \(S_0\) is
affine. If
\(x\in X(S)\) induces the unit section of \(X_{\mathscr J}\) and
normalizes \(Y\), then it centralizes \(Y\).}

This follows immediately by comparing the preceding corollary with
\SGARefLink{sga3:iii:corollary:2-14}{2.14}. Indeed, \SGARefLink{sga3:iii:corollary:4-33}{4.33} shows that the elements of \(C^0(Y_0,L_0^X)\)
which globally preserve \(Y\) are the elements of
\(H^0(Y_0,L_0^X)\), and \SGARefLink{sga3:iii:corollary:2-14}{2.14} shows that these are precisely the
elements which act trivially on the canonical immersion
\(Y\longrightarrow X\).

\medskip
\noindent\SGAInlineRefTarget{sga3:iii:paragraph:4-35}\textbf{4.35. ---}
Return to the general situation of \SGARefLink{sga3:iii:theorem:4-21}{4.21} and suppose that
\(Y_{\mathscr J}\) is smooth over \(S_{\mathscr J}\). By \SGARefLink{sga3:iii:lemma:4-25:item:ii}{4.25(ii)},
there is then an exact sequence of
\(Y_0\)-\(\mathscr O_{S_0}\)-modules
\[
 \tag{*}\label{sga3:iii:paragraph:4-35:equation:star}
 0\longrightarrow\operatorname{Lie}(Y_0/S_0)
 \longrightarrow\operatorname{Lie}(X_0/S_0)
 \longrightarrow\mathbf n_{Y_0/X_0}^{\vee}\longrightarrow0,
\]
and these are finite locally free \(\mathscr O_{S_0}\)-modules.

On the other hand, by SGA~1, II, \SGARefLink{sga3:ii:definition:4-10}{4.10}, every subscheme \(Y\) of
\(X\) lifting \(Y_{\mathscr J}\) and flat over \(S\) is smooth over
\(S\).\footnote{\label{fnref:III-132}Editor's note: What follows, together with
\SGARefLink{sga3:iii:proposition:4-35-1}{Proposition 4.35.1}, has been added; it is implicit in the original.
Cf. 4.38(5).} Suppose in addition that \(S_0\) and
\(Y_{\mathscr J}\) are affine. Since
\(\mathbf n_{Y_0/X_0}^{\vee}\) is a locally free
\(\mathscr O_{S_0}\)-module, the sequence \((*)\) remains exact after
tensoring with \(\mathscr J\) over \(\mathscr O_{S_0}\), and then
pulling back to \(Y_0^n\). Since the \(Y_0^n\) are affine, this gives
an exact sequence of complexes of abelian groups
\[
 0\longrightarrow C^*(Y_0,L_0^Y)
 \longrightarrow C^*(Y_0,L_0^X)
 \xrightarrow{d^*}C^*(Y_0,N_0)\longrightarrow0,
\]
and, in particular, a commutative diagram with exact rows:

\begin{center}
 \SGAThreeNativeInput[0.94\textwidth]
   {native_diagrams/exp3/diagram_062_exp3_p072_435_cochain_diagram.tex}
\end{center}

Let now \(Y\) and \(Y'\) be two subgroups of \(X\) lifting
\(Y_{\mathscr J}\) and flat, hence smooth, over \(S\). Since
\(Y_{\mathscr J}\) is affine, \SGARefLink{sga3:iii:proposition:0-15}{0.15} shows that \(Y\) and \(Y'\) are
isomorphic as schemes lifting \(Y_{\mathscr J}\); that is, there is
an isomorphism of \(S\)-schemes
\[
 f:Y\xrightarrow{\sim}Y'
\]

inducing the identity on \(Y_{\mathscr J}\). On the one hand, by
\SGARefLink{sga3:iii:proposition:1-2-4}{1.2.4}, \(f\) defines an element \(a\in C^1(Y_0,L_0^X)\) such that
\[
 f(y)=a(y_0)y
\]
for every \(y\in Y(S')\), \(S'\to S\), and by \SGARefLink{sga3:iii:lemma:4-27:item:i}{4.27(i)},
\[
 d^1(a)=d(Y,Y').
\]
Moreover, because \(Y\) and \(Y'\) are subgroups of \(X\), the last
element belongs to \(Z^1(Y_0,N_0)\) (cf. \SGARefLink{sga3:iii:theorem:4-21}{4.21}). Thus
\(\partial a\) is an element of \(Z^2(Y_0,L_0^Y)\), whose image
\(\overline{\partial a}\) in \(H^2(Y_0,L_0^Y)\) depends only on the
class \(\overline{d(Y,Y')}\in H^1(Y_0,N_0)\). This is precisely the
definition of the connecting map
\[
 \partial^1:H^1(Y_0,N_0)\longrightarrow H^2(Y_0,L_0^Y),
\]
so that
\[
 \partial^1\bigl(\overline{d(Y,Y')}\bigr)
 =\overline{\partial a}.
\]

On the other hand, transport the group structure of \(Y'\) by \(f\),
and let \(Y_1\) be the resulting group, whose underlying scheme is
therefore \(Y\). Thus its group law \(\mu_1\) is defined, for every
\(S'\to S\) and \(x,y\in Y(S')\), by
\[
 \mu_1(x,y)=f^{-1}\bigl(f(x)f(y)\bigr).
\]
By \SGARefLink{sga3:iii:remark:3-5-1}{3.5.1}, \(Y_1\) defines a cocycle
\[
 \delta(Y,Y_1)\in
 Z^2\bigl(Y_0,\operatorname{Lie}(Y_0/S_0)\bigr)
\]
such that, for every \(S'\to S\) and \(x,y\in Y(S')\),
\[
 \delta(Y,Y_1)(x_0,y_0)xy
 =\mu_1(x,y)=f^{-1}\bigl(f(x)f(y)\bigr).
\]
Put \(b=\delta(Y,Y_1)\). For every \(S'\to S\) and
\(x,y\in Y(S')\), we have
\((b(x_0,y_0)xy)_0=x_0y_0\). Hence
\(f(b(x_0,y_0)xy)\) is equal, on the one hand, to
\[
 a(x_0y_0)b(x_0,y_0)xy,
\]
and, on the other hand, to
\[
 f(x)f(y)=a(x_0)x\,a(y_0)y
 =a(x_0)\operatorname{Ad}(x_0)\bigl(a(y_0)\bigr)xy.
\]
Comparing these two expressions, we find that \(b(x_0,y_0)\) equals
\[
 \begin{aligned}
 &a(x_0y_0)^{-1}a(x_0)
   \operatorname{Ad}(x_0)\bigl(a(y_0)\bigr)\\
 &\qquad={}
 \operatorname{Ad}(x_0)\bigl(a(y_0)\bigr)-a(x_0y_0)+a(x_0)
 =(\partial a)(x_0,y_0).
 \end{aligned}
\]
Thus \(\delta(Y,Y_1)=\partial a\), and we have proved:

\medskip
\noindent\SGAInlineRefTarget{sga3:iii:proposition:4-35-1}\textbf{Proposition 4.35.1.}\textsuperscript{\normalfont\ref{fnref:III-132}}
\textit{Under the hypotheses of \SGARefLink{sga3:iii:theorem:4-21}{4.21}, suppose in addition that
\(S_0\) is affine and that \(Y_{\mathscr J}\) is smooth over
\(S_{\mathscr J}\) and affine. Let \(Y,Y'\) be two subgroups of
\(X\) lifting \(Y_{\mathscr J}\) and flat (hence smooth) over
\(S\). Let \(f:Y\xrightarrow{\sim}Y'\) be an isomorphism of
\(S\)-schemes inducing the identity on \(Y_{\mathscr J}\), and let
\(Y_1\) denote the group obtained by transporting the group structure
of \(Y'\) by \(f\). Then}
\[
 \partial^1\bigl(\overline{d(Y,Y')}\bigr)
 =\overline{\delta(Y,Y_1)}.
\]

